\documentclass[12pt,a4paper]{article}
\usepackage[utf8]{inputenc}
\usepackage[T1]{fontenc}
\usepackage[french]{babel}
\usepackage{geometry}
\usepackage{fancyhdr}
\usepackage{amsmath}
\usepackage{listings}
\usepackage{xcolor}
\usepackage{hyperref}
\usepackage{graphicx}
\usepackage{enumitem}
\usepackage{tcolorbox}

\geometry{margin=2.5cm}
\setlength{\headheight}{14.5pt}

\pagestyle{fancy}
\fancyhf{}
\fancyhead[L]{Résumé Pratique - Administration Unix}
\fancyhead[R]{\thepage}
\fancyfoot[C]{}

\lstset{
    language=bash,
    basicstyle=\ttfamily\small,
    keywordstyle=\color{blue},
    commentstyle=\color{gray},
    stringstyle=\color{red},
    numbers=left,
    numberstyle=\tiny\color{gray},
    breaklines=true,
    frame=single,
    showstringspaces=false
}

\title{\textbf{Résumé Pratique - Administration Unix/Linux}\\
\large Commandes et Explications de Cours - Examen Final}
\author{}
\date{}

\begin{document}

\maketitle
\thispagestyle{fancy}

\section*{Introduction}
Ce document présente un résumé pratique des commandes et explications de cours pour l'examen d'Administration Unix. Il couvre toutes les parties au programme :
\begin{itemize}
    \item Hiérarchie des systèmes de fichiers (TP1-0, TP1-1)
    \item Types de fichiers Linux
    \item Démarrage et arrêt du système, GRUB (TP2)
    \item Gestion des services (TP3)
    \item Gestion des paquetages (TP4)
    \item Gestion des utilisateurs et groupes (TP6)
    \item Gestion des disques : partitions, formatage, montage (TP7)
    \item Gestion de l'espace swap
\end{itemize}

\textbf{Note importante} : Les commandes nécessitant des privilèges root sont préfixées par \texttt{sudo} ou doivent être exécutées en tant que root.

\newpage
\tableofcontents
\newpage

% ============================================================
% SECTION 1 : HIÉRARCHIE DES SYSTÈMES DE FICHIERS (FHS)
% ============================================================

\section{Hiérarchie des Systèmes de Fichiers (FHS)}

\subsection{Définition du FHS}

Le \textbf{FHS (Filesystem Hierarchy Standard)} définit l'organisation standard des répertoires dans les systèmes Linux/Unix. Il garantit la cohérence entre les distributions et facilite l'administration.

\subsection{Commandes de base (TP1-0)}

\subsubsection{Navigation et exploration}
\begin{itemize}
    \item \texttt{pwd} : Afficher le répertoire courant (Print Working Directory)
    \item \texttt{cd [répertoire]} : Changer de répertoire
    \begin{itemize}
        \item \texttt{cd} ou \texttt{cd \textasciitilde} : Aller au répertoire personnel
        \item \texttt{cd ..} : Remonter d'un niveau
        \item \texttt{cd -} : Retourner au répertoire précédent
        \item \texttt{cd /} : Aller à la racine
    \end{itemize}
    \item \texttt{ls [options] [fichier/répertoire]} : Lister les fichiers
    \begin{itemize}
        \item \texttt{ls -l} : Format long (détails)
        \item \texttt{ls -a} : Afficher les fichiers cachés (commençant par .)
        \item \texttt{ls -h} : Tailles en format lisible (human-readable)
        \item \texttt{ls -R} : Récursif (sous-répertoires)
        \item \texttt{ls -t} : Trier par date de modification
        \item \texttt{ls -S} : Trier par taille
    \end{itemize}
\end{itemize}

\subsubsection{Manipulation de fichiers}
\begin{itemize}
    \item \texttt{cp [options] source destination} : Copier
    \begin{itemize}
        \item \texttt{cp -r} : Copier récursivement (répertoires)
        \item \texttt{cp -p} : Préserver les attributs (permissions, dates)
        \item \texttt{cp -i} : Mode interactif (demander confirmation)
    \end{itemize}
    \item \texttt{mv source destination} : Déplacer ou renommer
    \item \texttt{rm [options] fichier} : Supprimer
    \begin{itemize}
        \item \texttt{rm -r} : Supprimer récursivement
        \item \texttt{rm -f} : Forcer (sans confirmation)
        \item \texttt{rm -i} : Mode interactif
    \end{itemize}
    \item \texttt{mkdir [options] répertoire} : Créer un répertoire
    \begin{itemize}
        \item \texttt{mkdir -p} : Créer les répertoires parents si nécessaire
    \end{itemize}
    \item \texttt{rmdir répertoire} : Supprimer un répertoire vide
    \item \texttt{touch fichier} : Créer un fichier vide ou mettre à jour la date
\end{itemize}

\subsubsection{Affichage et édition}
\begin{itemize}
    \item \texttt{cat fichier} : Afficher le contenu complet
    \item \texttt{less fichier} ou \texttt{more fichier} : Afficher page par page
    \item \texttt{head [-n nombre] fichier} : Afficher les premières lignes (défaut: 10)
    \item \texttt{tail [-n nombre] fichier} : Afficher les dernières lignes (défaut: 10)
    \item \texttt{tail -f fichier} : Suivre les nouvelles lignes (logs en temps réel)
    \item \texttt{vi/vim fichier} : Éditeur de texte (mode insertion: i, sortie: Esc, sauvegarder: :wq)
    \item \texttt{nano fichier} : Éditeur de texte simple
\end{itemize}

\subsection{Répertoires principaux du FHS (TP1-1)}

\subsubsection{Répertoires système essentiels}
\begin{itemize}
    \item \textbf{/bin} : Exécutables essentiels pour tous les utilisateurs
    \begin{itemize}
        \item Commandes de base : ls, cp, mv, rm, cat, bash
        \item Disponibles même en mode single-user
    \end{itemize}
    \item \textbf{/sbin} : Exécutables système pour l'administration
    \begin{itemize}
        \item Nécessitent généralement root : shutdown, ifconfig, fsck, mount
    \end{itemize}
    \item \textbf{/lib} : Bibliothèques partagées (.so)
    \begin{itemize}
        \item Bibliothèques utilisées par les programmes de /bin et /sbin
    \end{itemize}
    \item \textbf{/usr/bin} : Exécutables utilisateur additionnels
    \item \textbf{/usr/sbin} : Exécutables système additionnels
    \item \textbf{/usr/lib} : Bibliothèques pour /usr/bin et /usr/sbin
\end{itemize}

\subsubsection{Répertoires de configuration et données}
\begin{itemize}
    \item \textbf{/etc} : Fichiers de configuration système
    \begin{itemize}
        \item \texttt{/etc/passwd} : Comptes utilisateurs
        \item \texttt{/etc/shadow} : Mots de passe cryptés
        \item \texttt{/etc/group} : Groupes
        \item \texttt{/etc/fstab} : Systèmes de fichiers à monter
        \item \texttt{/etc/hosts} : Résolution de noms locale
        \item \texttt{/etc/sudoers} : Configuration sudo
    \end{itemize}
    \item \textbf{/var} : Données variables (taille change)
    \begin{itemize}
        \item \texttt{/var/log} : Journaux système
        \item \texttt{/var/spool} : Files d'attente (mail, cron)
        \item \texttt{/var/tmp} : Fichiers temporaires persistants
    \end{itemize}
    \item \textbf{/tmp} : Fichiers temporaires (effacés au redémarrage)
    \item \textbf{/home} : Répertoires personnels des utilisateurs
    \item \textbf{/root} : Répertoire personnel de root
\end{itemize}

\subsubsection{Répertoires système}
\begin{itemize}
    \item \textbf{/dev} : Fichiers spéciaux (périphériques)
    \begin{itemize}
        \item \texttt{/dev/sda}, \texttt{/dev/sdb} : Disques SATA/SCSI
        \item \texttt{/dev/null} : "Poubelle" (ignore tout)
        \item \texttt{/dev/zero} : Source de zéros
    \end{itemize}
    \item \textbf{/proc} : Système de fichiers virtuel (interface avec le noyau)
    \begin{itemize}
        \item \texttt{/proc/cpuinfo} : Informations CPU
        \item \texttt{/proc/meminfo} : Informations mémoire
        \item \texttt{/proc/PID/} : Informations sur un processus
    \end{itemize}
    \item \textbf{/sys} : Système de fichiers virtuel (appareils et pilotes)
    \item \textbf{/boot} : Fichiers de démarrage (noyau, GRUB)
    \item \textbf{/mnt} : Point de montage temporaire
    \item \textbf{/media} : Points de montage pour médias amovibles
\end{itemize}

\subsubsection{Commandes d'exploration}
\begin{itemize}
    \item \texttt{find [chemin] [critères]} : Rechercher des fichiers
    \begin{itemize}
        \item \texttt{find . -name "*.txt"} : Par nom
        \item \texttt{find / -type f -size +100M} : Fichiers > 100 Mo
        \item \texttt{find . -mtime -7} : Modifiés dans les 7 derniers jours
    \end{itemize}
    \item \texttt{locate fichier} : Recherche rapide (base de données)
    \begin{itemize}
        \item \texttt{sudo updatedb} : Mettre à jour la base (root)
    \end{itemize}
    \item \texttt{which commande} : Chemin de l'exécutable
    \item \texttt{whereis commande} : Localiser binaire, source, man
    \item \texttt{type commande} : Type de commande (builtin, alias, fichier)
\end{itemize}

% ============================================================
% SECTION 2 : TYPES DE FICHIERS LINUX
% ============================================================

\section{Types de Fichiers Linux}

\subsection{Identification des types}

\begin{itemize}
    \item \texttt{file fichier} : Identifier le type d'un fichier
    \item \texttt{ls -l} : Afficher le type (premier caractère)
    \begin{itemize}
        \item \texttt{-} : Fichier régulier
        \item \texttt{d} : Répertoire
        \item \texttt{l} : Lien symbolique
        \item \texttt{c} : Périphérique caractère
        \item \texttt{b} : Périphérique bloc
        \item \texttt{p} : Pipe nommé (FIFO)
        \item \texttt{s} : Socket
    \end{itemize}
\end{itemize}

\subsection{Liens}

\subsubsection{Liens symboliques (soft links)}
\begin{itemize}
    \item \texttt{ln -s cible lien} : Créer un lien symbolique
    \begin{itemize}
        \item Pointe vers le nom du fichier
        \item Peut traverser les systèmes de fichiers
        \item Devient "cassé" si la cible est supprimée
    \end{itemize}
    \item Exemple : \texttt{ln -s /usr/bin/python3 /usr/local/bin/python}
\end{itemize}

\subsubsection{Liens physiques (hard links)}
\begin{itemize}
    \item \texttt{ln cible lien} : Créer un lien physique
    \begin{itemize}
        \item Partage le même inode
        \item Doit rester dans le même système de fichiers
        \item Le fichier n'est supprimé que quand tous les liens sont supprimés
    \end{itemize}
\end{itemize}

\subsection{Inodes}

\begin{itemize}
    \item \texttt{ls -i fichier} : Afficher le numéro d'inode
    \item \texttt{df -i} : Espace inode disponible
    \item \textbf{Important} : Le nom du fichier n'est PAS dans l'inode, mais dans le répertoire parent
\end{itemize}

% ============================================================
% SECTION 3 : DÉMARRAGE ET ARRÊT DU SYSTÈME, GRUB
% ============================================================

\section{Démarrage et Arrêt du Système, GRUB (TP2)}

\subsection{Processus de démarrage}

Séquence de démarrage :
\begin{enumerate}
    \item \textbf{Power on} : Alimentation
    \item \textbf{BIOS/UEFI} : Vérification matérielle
    \item \textbf{BootLoader (GRUB)} : Chargement du noyau
    \item \textbf{Linux Kernel} : Initialisation du noyau
    \item \textbf{Init/Systemd} : Processus PID 1
    \item \textbf{System Ready} : Système prêt
\end{enumerate}

\subsection{GRUB2 (GRand Unified Bootloader)}

\subsubsection{Fichiers GRUB2}
\begin{itemize}
    \item \texttt{/etc/default/grub} : Paramètres de configuration GRUB
    \item \texttt{/etc/grub.d/} : Scripts de génération du menu
    \item \texttt{/boot/grub2/grub.cfg} : Fichier de configuration final (généré automatiquement)
    \begin{itemize}
        \item \textbf{Ne jamais éditer directement !}
    \end{itemize}
\end{itemize}

\subsubsection{Commandes GRUB2}
\begin{itemize}
    \item \texttt{sudo grub2-mkconfig -o /boot/grub2/grub.cfg} : Régénérer grub.cfg
    \begin{itemize}
        \item À exécuter après modification de /etc/default/grub
    \end{itemize}
    \item \texttt{sudo grub2-install /dev/sda} : Installer GRUB sur un disque
    \item \texttt{sudo grub2-set-default 0} : Définir l'entrée par défaut
\end{itemize}

\subsubsection{Configuration /etc/default/grub}
\begin{lstlisting}
GRUB_TIMEOUT=5              # Délai avant démarrage automatique
GRUB_DEFAULT=saved          # Entrée par défaut
GRUB_CMDLINE_LINUX="..."    # Options du noyau
GRUB_DISABLE_SUBMENU=true   # Désactiver sous-menus
\end{lstlisting}

\subsubsection{Initramfs (Initial RAM Filesystem)}
\begin{itemize}
    \item Système de fichiers temporaire en RAM
    \item Contient les modules nécessaires pour accéder au disque racine
    \item Nécessaire si racine sur LVM, RAID, ou système crypté
    \item \texttt{sudo dracut --force} : Régénérer l'initramfs
\end{itemize}

\subsection{Arrêt et redémarrage}

\subsubsection{Commandes systemd}
\begin{itemize}
    \item \texttt{sudo systemctl poweroff} : Arrêt et hors-tension
    \item \texttt{sudo systemctl halt} : Arrêt sans hors-tension
    \item \texttt{sudo systemctl reboot} : Redémarrage
    \item \texttt{sudo systemctl suspend} : Suspension (RAM)
    \item \texttt{sudo systemctl hibernate} : Hibernation (disque)
\end{itemize}

\subsubsection{Commande shutdown}
\begin{itemize}
    \item \texttt{sudo shutdown [OPTION] [HEURE] [MESSAGE]}
    \begin{itemize}
        \item \texttt{sudo shutdown -h now} : Arrêt immédiat
        \item \texttt{sudo shutdown -r +10} : Redémarrage dans 10 minutes
        \item \texttt{sudo shutdown -c} : Annuler un arrêt programmé
    \end{itemize}
\end{itemize}

% ============================================================
% SECTION 4 : GESTION DES SERVICES (TP3)
% ============================================================

\section{Gestion des Services (TP3)}

\subsection{Systemd - Gestionnaire de services moderne}

\subsubsection{Concepts}
\begin{itemize}
    \item \textbf{Unités} : Services, targets, mounts, sockets, timers
    \item \textbf{Dépendances} : Démarrage ordonné selon dépendances
    \item \textbf{Parallélisation} : Démarrage parallèle quand possible
    \item \textbf{Journalisation} : Logs intégrés (journalctl)
\end{itemize}

\subsubsection{Commandes de contrôle des services}
\begin{itemize}
    \item \texttt{sudo systemctl start service} : Démarrer un service
    \item \texttt{sudo systemctl stop service} : Arrêter un service
    \item \texttt{sudo systemctl restart service} : Redémarrer un service
    \item \texttt{sudo systemctl reload service} : Recharger la configuration
    \item \texttt{systemctl status service} : État détaillé d'un service
    \item \texttt{systemctl is-active service} : Vérifier si actif
    \item \texttt{systemctl is-enabled service} : Vérifier activation au boot
\end{itemize}

\subsubsection{Activation au démarrage}
\begin{itemize}
    \item \texttt{sudo systemctl enable service} : Activer au démarrage
    \item \texttt{sudo systemctl disable service} : Désactiver au démarrage
    \item \texttt{sudo systemctl mask service} : Bloquer (empêcher démarrage)
    \item \texttt{sudo systemctl unmask service} : Débloquer
\end{itemize}

\subsubsection{Listes et informations}
\begin{itemize}
    \item \texttt{systemctl list-units} : Lister toutes les unités
    \item \texttt{systemctl list-units --type=service} : Lister les services
    \item \texttt{systemctl list-unit-files} : Lister les fichiers d'unités
    \item \texttt{systemctl list-dependencies service} : Dépendances d'un service
\end{itemize}

\subsection{Targets (équivalents des runlevels)}

\begin{itemize}
    \item \texttt{poweroff.target} : Arrêt (runlevel 0)
    \item \texttt{rescue.target} : Mode maintenance (runlevel 1)
    \item \texttt{multi-user.target} : Multi-utilisateurs texte (runlevel 3)
    \item \texttt{graphical.target} : Multi-utilisateurs graphique (runlevel 5)
    \item \texttt{reboot.target} : Redémarrage (runlevel 6)
\end{itemize}

\subsubsection{Commandes targets}
\begin{itemize}
    \item \texttt{systemctl get-default} : Target par défaut
    \item \texttt{sudo systemctl set-default target} : Définir target par défaut
    \item \texttt{sudo systemctl isolate target} : Changer de target
\end{itemize}

\subsection{Journalisation avec journalctl}

\subsubsection{Commandes journalctl}
\begin{itemize}
    \item \texttt{journalctl} : Afficher tous les logs
    \item \texttt{journalctl -f} : Suivre les nouveaux logs (comme tail -f)
    \item \texttt{journalctl -b} : Logs du démarrage actuel
    \item \texttt{journalctl -b -1} : Logs du démarrage précédent
    \item \texttt{journalctl -u service} : Logs d'un service spécifique
    \item \texttt{journalctl -p err} : Logs d'erreur et supérieurs
    \item \texttt{journalctl --since "1 hour ago"} : Logs depuis 1 heure
    \item \texttt{journalctl --until "2024-01-01 12:00"} : Logs jusqu'à une date
    \item \texttt{journalctl -k} : Logs du noyau uniquement
\end{itemize}

\subsubsection{Options de formatage}
\begin{itemize}
    \item \texttt{journalctl -n 50} : 50 dernières lignes
    \item \texttt{journalctl --no-pager} : Sans pagination
    \item \texttt{journalctl -o json} : Format JSON
\end{itemize}

% ============================================================
% SECTION 5 : GESTION DES PAQUETAGES (TP4)
% ============================================================

\section{Gestion des Paquetages (TP4)}

\subsection{Concepts}

\begin{itemize}
    \item \textbf{Paquet} : Archive contenant logiciel + métadonnées + scripts
    \item \textbf{Dépendances} : Paquets requis pour fonctionner
    \item \textbf{Dépôts (repositories)} : Sources de paquets (réseau ou local)
\end{itemize}

\subsection{Format RPM (RedHat Package Manager)}

\subsubsection{Commandes RPM (bas niveau)}
\begin{itemize}
    \item \texttt{sudo rpm -i package.rpm} : Installer un paquet
    \item \texttt{sudo rpm -e package} : Désinstaller (erase)
    \item \texttt{sudo rpm -U package.rpm} : Mettre à jour (upgrade)
    \item \texttt{rpm -qa} : Lister tous les paquets installés
    \item \texttt{rpm -q package} : Version installée
    \item \texttt{rpm -qi package} : Informations sur un paquet
    \item \texttt{rpm -ql package} : Liste des fichiers installés
    \item \texttt{rpm -qf fichier} : Paquet auquel appartient un fichier
    \item \texttt{rpm -V package} : Vérifier intégrité
    \item \texttt{rpm -K package.rpm} : Vérifier signature
\end{itemize}

\subsection{Gestionnaire YUM (Yellowdog Updater Modified)}

\subsubsection{Commandes YUM}
\begin{itemize}
    \item \texttt{yum list all} : Lister tous les paquets disponibles
    \item \texttt{yum list installed} : Lister les paquets installés
    \item \texttt{yum list available} : Lister les paquets disponibles
    \item \texttt{yum list updates} : Lister les mises à jour disponibles
    \item \texttt{yum info package} : Informations sur un paquet
    \item \texttt{yum search mot\_clé} : Rechercher des paquets
    \item \texttt{sudo yum install package} : Installer
    \item \texttt{sudo yum update [package]} : Mettre à jour (tout ou un paquet)
    \item \texttt{sudo yum remove package} : Désinstaller
    \item \texttt{yum check-update} : Vérifier mises à jour disponibles
    \item \texttt{yum provides fichier} : Chercher paquet contenant un fichier
    \item \texttt{sudo yum clean all} : Vider le cache
    \item \texttt{sudo yum makecache} : Mettre à jour la liste des paquets
\end{itemize}

\subsubsection{Gestion des groupes}
\begin{itemize}
    \item \texttt{yum grouplist} : Lister les groupes de paquets
    \item \texttt{sudo yum groupinstall "nom\_groupe"} : Installer un groupe
    \item \texttt{sudo yum groupremove "nom\_groupe"} : Désinstaller un groupe
    \item \texttt{yum groupinfo "nom\_groupe"} : Informations sur un groupe
\end{itemize}

\subsubsection{Historique}
\begin{itemize}
    \item \texttt{yum history} : Historique des transactions
    \item \texttt{yum history info ID} : Détails d'une transaction
    \item \texttt{sudo yum history undo ID} : Annuler une transaction
\end{itemize}

\subsection{Gestionnaire DNF (Dandified Yum)}

DNF est le successeur de YUM (Fedora 22+, RHEL 8+). Syntaxe identique :
\begin{itemize}
    \item \texttt{sudo dnf install/remove/update} : Gestion des paquets
    \item \texttt{dnf search/info/provides} : Recherche et informations
    \item \texttt{dnf list [available/installed/updates]} : Listes
    \item \texttt{sudo dnf groupinstall/groupremove} : Gestion des groupes
    \item \texttt{dnf history} : Historique
\end{itemize}

\subsection{Dépôts de paquetages}

\subsubsection{Configuration}
\begin{itemize}
    \item Fichiers \texttt{.repo} dans \texttt{/etc/yum.repos.d/}
    \item \texttt{yum repolist all} : Lister tous les dépôts
    \item \texttt{yum repolist} : Lister les dépôts actifs
    \item \texttt{sudo yum-config-manager --enable nom\_dépôt} : Activer un dépôt
    \item \texttt{sudo yum-config-manager --disable nom\_dépôt} : Désactiver
\end{itemize}

% ============================================================
% SECTION 6 : GESTION DES UTILISATEURS ET GROUPES (TP6)
% ============================================================

\section{Gestion des Utilisateurs et Groupes (TP6)}

\subsection{Fichiers de configuration}

\subsubsection{/etc/passwd}
Format : \texttt{login:passwd:uid:gid:comment:home:shell}
\begin{itemize}
    \item \texttt{login} : Nom de connexion
    \item \texttt{passwd} : x (indique présence de /etc/shadow)
    \item \texttt{uid} : User ID (0 = root, < 1000 = système, >= 1000 = utilisateurs)
    \item \texttt{gid} : Group ID primaire
    \item \texttt{comment} : Nom complet (GECOS)
    \item \texttt{home} : Répertoire personnel
    \item \texttt{shell} : Shell par défaut (/bin/bash, /sbin/nologin)
\end{itemize}

\subsubsection{/etc/shadow}
Format : \texttt{login:passwd:last:may:must:warn:inactif:disable}
\begin{itemize}
    \item \texttt{passwd} : Mot de passe crypté
    \item \texttt{last} : Jours depuis dernière modification (depuis 1/1/1970)
    \item \texttt{may} : Jours minimum avant changement
    \item \texttt{must} : Jours maximum avant expiration
    \item \texttt{warn} : Jours d'avertissement avant expiration
    \item \texttt{inactif} : Période d'inactivité maximale
    \item \texttt{disable} : Date de fermeture du compte
\end{itemize}

\subsubsection{/etc/group}
Format : \texttt{nom:passwd:gid:membres}
\begin{itemize}
    \item Liste des membres séparés par des virgules
\end{itemize}

\subsection{Gestion des utilisateurs}

\subsubsection{Création}
\begin{itemize}
    \item \texttt{sudo useradd [options] nom\_utilisateur}
    \begin{itemize}
        \item \texttt{-u UID} : Spécifier UID
        \item \texttt{-c "commentaire"} : Nom complet
        \item \texttt{-d répertoire} : Répertoire personnel
        \item \texttt{-e AAAA-MM-JJ} : Date d'expiration
        \item \texttt{-g groupe} : Groupe primaire
        \item \texttt{-G groupe1,groupe2} : Groupes secondaires
        \item \texttt{-m} : Créer le répertoire personnel
        \item \texttt{-s shell} : Shell par défaut
        \item \texttt{-r} : Créer un compte système
    \end{itemize}
    \item \texttt{sudo passwd utilisateur} : Définir le mot de passe
\end{itemize}

\subsubsection{Modification}
\begin{itemize}
    \item \texttt{sudo usermod [options] utilisateur}
    \begin{itemize}
        \item \texttt{-s shell} : Changer le shell
        \item \texttt{-c "commentaire"} : Modifier le commentaire
        \item \texttt{-aG groupe} : Ajouter à des groupes (sans écraser)
        \item \texttt{-d -m répertoire} : Déplacer le home
        \item \texttt{-L} : Verrouiller le compte
        \item \texttt{-U} : Déverrouiller
        \item \texttt{-e date} : Expirer le compte
    \end{itemize}
\end{itemize}

\subsubsection{Suppression}
\begin{itemize}
    \item \texttt{sudo userdel utilisateur} : Supprimer (garde le home)
    \item \texttt{sudo userdel -r utilisateur} : Supprimer avec le home
    \item \texttt{sudo userdel -rf utilisateur} : Forcer même si connecté
\end{itemize}

\subsubsection{Gestion des mots de passe}
\begin{itemize}
    \item \texttt{passwd} : Changer son propre mot de passe
    \item \texttt{sudo passwd [options] utilisateur} : Gérer mot de passe utilisateur
    \begin{itemize}
        \item \texttt{-d} : Supprimer le mot de passe
        \item \texttt{-l} : Verrouiller
        \item \texttt{-u} : Déverrouiller
        \item \texttt{-e} : Faire expirer
        \item \texttt{-n jours} : Jours minimum avant changement
        \item \texttt{-x jours} : Jours maximum avant expiration
        \item \texttt{-w jours} : Délai d'avertissement
        \item \texttt{-i jours} : Délai avant désactivation
    \end{itemize}
    \item \texttt{sudo chage [options] utilisateur} : Modifier validité du mot de passe
    \begin{itemize}
        \item \texttt{-l} : Afficher informations
        \item \texttt{-m jours} : Jours minimum
        \item \texttt{-M jours} : Jours maximum
        \item \texttt{-W jours} : Jours d'avertissement
        \item \texttt{-I jours} : Jours d'inactivité
        \item \texttt{-E date} : Date d'expiration
    \end{itemize}
\end{itemize}

\subsection{Gestion des groupes}

\subsubsection{Création et modification}
\begin{itemize}
    \item \texttt{sudo groupadd [options] nom\_groupe}
    \begin{itemize}
        \item \texttt{-g GID} : Spécifier GID
        \item \texttt{-r} : Créer un groupe système
    \end{itemize}
    \item \texttt{sudo groupmod [options] nom\_groupe}
    \begin{itemize}
        \item \texttt{-g GID} : Changer GID
        \item \texttt{-n nouveau\_nom} : Renommer
    \end{itemize}
    \item \texttt{sudo groupdel nom\_groupe} : Supprimer
\end{itemize}

\subsubsection{Gestion des membres}
\begin{itemize}
    \item \texttt{sudo gpasswd -a utilisateur groupe} : Ajouter un membre
    \item \texttt{sudo gpasswd -d utilisateur groupe} : Retirer un membre
    \item \texttt{sudo gpasswd -M utilisateur1,utilisateur2 groupe} : Définir liste de membres
\end{itemize}

\subsection{Permissions et droits d'accès}

\subsubsection{Permissions traditionnelles}
\begin{itemize}
    \item Trois catégories : propriétaire (u), groupe (g), autres (o)
    \item Trois droits : lecture (r=4), écriture (w=2), exécution (x=1)
    \item Mode octal : somme des droits (ex: 755 = rwxr-xr-x)
\end{itemize}

\subsubsection{Modification des permissions}
\begin{itemize}
    \item \texttt{chmod [mode] fichier} : Modifier permissions
    \begin{itemize}
        \item Mode symbolique : \texttt{u+x}, \texttt{g-rw}, \texttt{o=r}, \texttt{a+r}
        \item Mode absolu : \texttt{755}, \texttt{644}, \texttt{600}, \texttt{777}
    \end{itemize}
    \item \texttt{sudo chmod -R mode répertoire} : Récursif
\end{itemize}

\subsubsection{Changement de propriétaire}
\begin{itemize}
    \item \texttt{sudo chown utilisateur[:groupe] fichier} : Changer propriétaire
    \item \texttt{sudo chgrp groupe fichier} : Changer groupe
    \item \texttt{sudo chown -R utilisateur:groupe répertoire} : Récursif
\end{itemize}

\subsection{Sudo - Gestion des privilèges}

\subsubsection{Configuration}
\begin{itemize}
    \item \texttt{/etc/sudoers} : Fichier de configuration
    \item \texttt{sudo visudo} : Éditer de façon sécurisée (vérifie syntaxe)
    \item Format : \texttt{utilisateur hôte=(utilisateur\_cible) commandes}
    \item Exemples :
    \begin{itemize}
        \item \texttt{jean ALL=(ALL:ALL) ALL} : Privilèges complets
        \item \texttt{\%wheel ALL=(ALL) ALL} : Groupe wheel avec privilèges
        \item \texttt{marie ALL=(root) NOPASSWD:/usr/sbin/reboot} : Commande spécifique sans mot de passe
    \end{itemize}
\end{itemize}

% ============================================================
% SECTION 7 : GESTION DES DISQUES (TP7)
% ============================================================

\section{Gestion des Disques : Partitions, Formatage et Montage (TP7)}

\subsection{Types de disques}

\begin{itemize}
    \item \textbf{HDD} : Disque dur mécanique, économique, lent
    \item \textbf{SSD} : Disque électronique, rapide, silencieux
    \item \textbf{NVMe} : Très haute performance (PCIe)
\end{itemize}

\subsection{Partitionnement}

\subsubsection{MBR vs GPT}
\begin{itemize}
    \item \textbf{MBR (Master Boot Record)}
    \begin{itemize}
        \item Limite : 4 partitions primaires maximum
        \item Taille max : 2.2 To par partition
        \item Compatible BIOS
    \end{itemize}
    \item \textbf{GPT (GUID Partition Table)}
    \begin{itemize}
        \item Jusqu'à 128 partitions
        \item Taille très grande (jusqu'à 9.4 Zo)
        \item Compatible UEFI et BIOS
    \end{itemize}
\end{itemize}

\subsubsection{Organisation des disques sous Linux}
\begin{itemize}
    \item Format : \texttt{/dev/sdXY} ou \texttt{/dev/nvme0n1pX}
    \begin{itemize}
        \item \texttt{/dev/sda} : Premier disque SATA/SCSI
        \item \texttt{/dev/sda1} : Première partition du premier disque
        \item \texttt{/dev/sdb} : Deuxième disque
    \end{itemize}
    \item Types de bus : hd (IDE), sd (SATA/SCSI/SSD), nvme (NVMe)
    \item UUID : Identifiant unique universel
\end{itemize}

\subsubsection{Commandes d'information}
\begin{itemize}
    \item \texttt{lsblk} : Afficher informations sur disques et partitions
    \item \texttt{lsblk -f} : Avec type de système de fichiers
    \item \texttt{sudo blkid} : Obtenir UUID
    \item \texttt{df -h} : Taux de remplissage (human-readable)
    \item \texttt{df -i} : Espace inode disponible
    \item \texttt{du -h [chemin]} : Espace occupé par une arborescence
    \item \texttt{du -sh répertoire} : Taille totale d'un répertoire
\end{itemize}

\subsubsection{Création de partitions}
\begin{itemize}
    \item \texttt{sudo fdisk /dev/sdX} : Manipuler table de partitions (interactif)
    \begin{itemize}
        \item \texttt{n} : Nouvelle partition
        \item \texttt{d} : Supprimer partition
        \item \texttt{t} : Changer type
        \item \texttt{p} : Afficher table
        \item \texttt{w} : Sauvegarder et quitter
        \item \texttt{q} : Quitter sans sauvegarder
        \item \texttt{g} : Créer nouvelle table GPT
    \end{itemize}
    \item \texttt{sudo gdisk /dev/sdX} : Pour GPT (similaire à fdisk)
    \item \texttt{sudo parted} : Mode interactif, script ou mixte
\end{itemize}

\subsection{Systèmes de fichiers}

\subsubsection{Types de systèmes de fichiers}
\begin{itemize}
    \item \textbf{ext2} : Stable et robuste (sans journalisation)
    \item \textbf{ext3} : ext2 + journalisation
    \item \textbf{ext4} : Volumes > 1To, timestamps nanoseconde, extents
    \item \textbf{XFS} : Performance fichiers volumineux
    \item \textbf{Btrfs} : Snapshots et compression
\end{itemize}

\subsubsection{Initialisation (formatage)}
\begin{itemize}
    \item \texttt{sudo mkfs -t type partition}
    \begin{itemize}
        \item \texttt{sudo mkfs.ext4 /dev/sda1} : Formater en ext4
        \item \texttt{sudo mkfs.xfs /dev/sda1} : Formater en XFS
    \end{itemize}
    \item \texttt{sudo lsblk -f} : Vérifier le type après formatage
\end{itemize}

\subsection{Montage d'une partition}

\subsubsection{Commandes de montage}
\begin{itemize}
    \item \texttt{sudo mount partition point\_de\_montage} : Monter
    \item \texttt{sudo mount -t type partition point} : Spécifier type
    \item \texttt{sudo mount -a} : Monter tous les SF de /etc/fstab
    \item \texttt{sudo umount point\_de\_montage} : Démontage
    \item \texttt{sudo umount -l point\_de\_montage} : Démontage paresseux (lazy)
\end{itemize}

\subsubsection{Fichier /etc/fstab}
Configuration statique des montages. Structure :
\begin{verbatim}
périphérique  point_montage  type  options  dump  fsck
\end{verbatim}

Exemple :
\begin{verbatim}
UUID=1234-5678  /mnt/data  ext4  defaults  0  2
/dev/sdb1       /mnt/usb   vfat  defaults  0  0
\end{verbatim}

Options courantes :
\begin{itemize}
    \item \texttt{defaults} : rw, suid, dev, exec, auto, nouser, async
    \item \texttt{ro} : Lecture seule
    \item \texttt{noexec} : Interdire exécution
    \item \texttt{nosuid} : Ignorer setuid/setgid
    \item \texttt{user} : Permettre montage par utilisateur
    \item \texttt{noauto} : Ne pas monter automatiquement
\end{itemize}

Champs :
\begin{itemize}
    \item \texttt{dump} : Activer sauvegarde (1 ou 0)
    \item \texttt{fsck} : Ordre de vérification (0=pas de vérif, 1=racine, 2+=autres)
\end{itemize}

% ============================================================
% SECTION 8 : GESTION DE L'ESPACE SWAP
% ============================================================

\section{Gestion de l'Espace Swap}

\subsection{Définition}

Zone de swap utilisée lorsque la mémoire physique (RAM) est remplie. Les pages mémoire inactives sont déplacées dans la zone de swap.

\subsection{Configuration d'un fichier swap}

\subsubsection{Étapes de configuration}
\begin{enumerate}
    \item Créer un fichier swap :
    \begin{itemize}
        \item \texttt{sudo fallocate -l 1G /swapfile} : Créer fichier 1 Go
        \item ou \texttt{sudo dd if=/dev/zero of=/swapfile bs=1M count=1024}
    \end{itemize}
    \item Verrouiller permissions : \texttt{sudo chmod 600 /swapfile}
    \item Marquer comme swap : \texttt{sudo mkswap /swapfile}
    \item Activer : \texttt{sudo swapon /swapfile}
    \item Rendre permanent : Ajouter dans /etc/fstab
    \begin{verbatim}
    /swapfile  none  swap  defaults  0  0
    \end{verbatim}
\end{enumerate}

\subsection{Commandes swap}

\begin{itemize}
    \item \texttt{swapon --show} : Afficher swaps actifs
    \item \texttt{free -h} : État mémoire et swap (human-readable)
    \item \texttt{free -m} : En mégaoctets
    \item \texttt{sudo swapoff /swapfile} : Désactiver
    \item \texttt{cat /proc/swaps} : Informations sur les swaps
\end{itemize}

\subsection{Partition swap}

Pour créer une partition swap :
\begin{enumerate}
    \item Créer une partition avec fdisk/gdisk
    \item Changer le type en "Linux swap" (82 pour MBR, 8200 pour GPT)
    \item \texttt{sudo mkswap /dev/sda2} : Formater comme swap
    \item \texttt{sudo swapon /dev/sda2} : Activer
    \item Ajouter dans /etc/fstab :
    \begin{verbatim}
    /dev/sda2  none  swap  defaults  0  0
    \end{verbatim}
\end{enumerate}

% ============================================================
% CONCLUSION
% ============================================================

\section*{Conclusion}

Ce résumé pratique couvre toutes les parties au programme de l'examen d'Administration Unix :
\begin{itemize}
    \item Hiérarchie des systèmes de fichiers et commandes de base
    \item Types de fichiers Linux
    \item Démarrage, arrêt et configuration GRUB
    \item Gestion des services avec systemd
    \item Gestion des paquetages (RPM, YUM, DNF)
    \item Gestion des utilisateurs, groupes et permissions
    \item Gestion des disques : partitionnement, formatage, montage
    \item Gestion de l'espace swap
\end{itemize}

\textbf{Rappel important} : La plupart des commandes d'administration nécessitent des privilèges root (préfixe \texttt{sudo} ou connexion en tant que root).

\end{document}

